% ===========================================================================
%  Chapitre 7 — Suites numériques
%  PE 2026, composante ALGÈBRE
% ===========================================================================
\bmtchapitre{Suites numériques}
  {Utiliser les suites numériques pour résoudre des problèmes : raisonnement
   par récurrence, convergence, suites du type $u_{n}=f(n)$ et
   $u_{n+1}=g(u_{n})$, suites arithmétiques, géométriques, stationnaires et
   adjacentes.}
  {Algèbre \textbullet\ environ 12 heures}

Une suite est une liste infinie de nombres, numérotée par les entiers. Cette
définition tient en une ligne et ne dit rien de l'intérêt de l'objet : ce qui
compte, c'est qu'une suite décrit un processus qui se répète — un capital qui
produit des intérêts, une dette qui s'amortit, un chiffre d'affaires qui
croît d'année en année.

Deux outils font la force de ce chapitre. Le raisonnement par récurrence,
d'abord, qui permet de démontrer une propriété pour une infinité d'entiers en
n'écrivant que deux lignes. L'inégalité des accroissements finis, ensuite, qui
transforme une suite définie par $u_{n+1} = f(u_{n})$ en une suite dont on
maîtrise la distance à sa limite.

% ---------------------------------------------------------------------------
\begin{revision}
Les acquis de première et les chapitres d'analyse suffisent.

\begin{enumerate}[label=\textbf{\arabic*.}, leftmargin=*, itemsep=0.35em]
  \item Une suite arithmétique de premier terme $3$ et de raison $2$ :
        donner $u_{5}$ et $u_{n}$.
  \item Une suite géométrique de premier terme $1$ et de raison $\frac12$ :
        donner $v_{4}$ et $v_{n}$.
  \item Calculer $1+2+3+\cdots+100$.
  \item Rappeler l'énoncé de l'inégalité des accroissements finis.
  \item Résoudre $\left(\frac23\right)^{n} \leq 10^{-3}$.
  \item Que vaut $\lim_{n\to+\infty}q^{n}$ selon les valeurs de $q$ ?
\end{enumerate}
\end{revision}

\begin{automatismes}
Sans calculatrice, sans brouillon.

\begin{enumerate}[label=\textbf{\alph*)}, leftmargin=*, itemsep=0.25em]
  \item $u_{n} = 3n-1$ : donner $u_{0}$ et $u_{10}$
  \item $v_{n} = 2^{n}$ : donner $v_{0}$ et $v_{5}$
  \item Somme $1+2+\cdots+n$
  \item Somme $1+q+q^{2}+\cdots+q^{n}$ pour $q \neq 1$
  \item $\lim\left(\frac12\right)^{n}$
  \item $\lim\left(\frac32\right)^{n}$
  \item $\lim (-1)^{n}$
  \item $\lim \frac{1}{n}$
\end{enumerate}
\end{automatismes}

\begin{corrige}[automatismes]
a) $-1$ et $29$ \quad b) $1$ et $32$ \quad c) $\frac{n(n+1)}{2}$ \quad
d) $\frac{1-q^{n+1}}{1-q}$ \quad e) $0$ \quad f) $+\infty$ \quad
g) pas de limite \quad h) $0$.
\end{corrige}

% ===========================================================================
\section{Définition et modes de génération}

\begin{definition}[suite numérique]
Une \textbf{suite numérique} est une fonction de $\N$ (ou d'une partie de
$\N$) dans $\R$. L'image de l'entier $n$ se note $u_{n}$ et s'appelle le
\textbf{terme de rang} $n$. La suite elle-même se note $(u_{n})$.
\end{definition}

\begin{propriete}[deux modes de génération]
Une suite est donnée soit \textbf{explicitement}, par une formule
$u_{n} = f(n)$, soit \textbf{par récurrence}, par un premier terme et une
relation $u_{n+1} = g(u_{n})$.
\end{propriete}

\begin{visualisation}[la construction en escalier]
\begin{center}
\begin{tikzpicture}
\begin{axis}[
  width = 9.6cm, height = 6.4cm,
  axis lines = middle, xlabel = {$x$},
  xmin = -0.3, xmax = 5.4, ymin = -0.3, ymax = 4.3,
  xtick = {1, 1.75, 2.125}, xticklabels = {$u_{0}$, $u_{1}$, $u_{2}$},
  ytick = \empty,
  axis line style = {bmtGris},
  tick label style = {font=\footnotesize, color=bmtGris}, clip = true,
]
  \addplot[bmtLaterite, line width=1.2pt, domain=0:4, samples=2] {x};
  \addplot[bmtIndigo, line width=1.3pt, domain=0:4, samples=120]
    {0.5*x + 1.25};
  \draw[bmtVetiver, line width=0.9pt] (axis cs:1,0) -- (axis cs:1,1.75)
    -- (axis cs:1.75,1.75) -- (axis cs:1.75,2.125)
    -- (axis cs:2.125,2.125) -- (axis cs:2.125,2.3125)
    -- (axis cs:2.3125,2.3125) -- (axis cs:2.3125,2.406);
  \filldraw[bmtVetiver] (axis cs:1,1.75) circle (1.6pt);
  \filldraw[bmtVetiver] (axis cs:1.75,2.125) circle (1.6pt);
  \filldraw[bmtLaterite] (axis cs:2.5,2.5) circle (2pt);
  \node[bmtLaterite, font=\footnotesize, anchor=north west]
    at (axis cs:2.58,2.38) {$\ell$};
  \node[bmtIndigo, font=\footnotesize, anchor=west]
    at (axis cs:4.08,3.25) {$y = g(x)$};
  \node[bmtLaterite, font=\footnotesize, anchor=south east]
    at (axis cs:3.9,3.95) {$y = x$};
\end{axis}
\end{tikzpicture}
\end{center}

Pour construire les termes d'une suite définie par $u_{n+1} = g(u_{n})$, on
part de $u_{0}$ sur l'axe des abscisses, on monte jusqu'à la courbe, puis on
se déplace horizontalement jusqu'à la droite $y = x$. La figure montre
aussitôt ce qu'aucun calcul ne montre : la suite converge vers le point
d'intersection de la courbe et de la droite.
\end{visualisation}

% ===========================================================================
\section{Raisonnement par récurrence}

\begin{theoreme}[principe de récurrence]
Soit $P(n)$ une propriété portant sur l'entier $n$, et $n_{0}$ un entier. Si
$P(n_{0})$ est vraie, et si pour tout $n \geq n_{0}$, $P(n)$ vraie entraîne
$P(n+1)$ vraie, alors $P(n)$ est vraie pour tout entier $n \geq n_{0}$.
\end{theoreme}

\begin{methode}{rédiger une récurrence}
\begin{enumerate}[label=\textbf{\arabic*.}, leftmargin=*, itemsep=0.15em]
  \item \emph{Énoncé.} Écrire explicitement la propriété $P(n)$.
  \item \emph{Initialisation.} Vérifier $P(n_{0})$ par le calcul.
  \item \emph{Hérédité.} Supposer $P(n)$ vraie, démontrer $P(n+1)$.
  \item \emph{Conclusion.} Invoquer le principe de récurrence.
\end{enumerate}
\end{methode}

\begin{avertissement}
L'hypothèse de récurrence doit être \emph{utilisée} dans l'hérédité. Une
démonstration qui n'y fait jamais appel n'est pas une récurrence.
\end{avertissement}

\begin{exemple}[une récurrence complète]
Montrons que pour tout $n \in \N$, $4^{n}+6n-1$ est divisible par $9$.

\emph{Initialisation.} Pour $n = 0$, $4^{0}+0-1 = 0$, divisible par $9$.

\emph{Hérédité.} Supposons $4^{n}+6n-1 = 9k$ avec $k$ entier. Alors
\[
  4^{n+1}+6(n+1)-1 = 4\left(4^{n}+6n-1\right)-18n+9
  = 9\left(4k-2n+1\right),
\]
qui est bien divisible par $9$.

\emph{Conclusion.} La propriété est vraie pour tout entier naturel $n$.
\end{exemple}

\begin{entrainement}[récurrence]
\exo[\niveau{1}] Démontrer que $\displaystyle\sum_{k=1}^{n}k
= \frac{n(n+1)}{2}$.

\exo[\niveau{2}] Démontrer que pour tout $n \geq 1$, $2^{n} \geq n+1$.

\exo[\niveau{3}] Soit $(u_{n})$ définie par $u_{0} = 1$ et
$u_{n+1} = \sqrt{2+u_{n}}$. Démontrer que pour tout $n$,
$1 \leq u_{n} \leq 2$.
\end{entrainement}

% ===========================================================================
\section{Convergence et divergence}

\begin{definition}[limite d'une suite]
La suite $(u_{n})$ \textbf{converge} vers le réel $\ell$ si tout intervalle
ouvert contenant $\ell$ contient tous les termes de la suite à partir d'un
certain rang. Une suite qui ne converge pas est dite \textbf{divergente}.
\end{definition}

\begin{theoreme}[convergence monotone]
Toute suite croissante et majorée converge. Toute suite décroissante et
minorée converge.
\end{theoreme}

\begin{avertissement}
Ce théorème affirme l'existence de la limite, jamais sa valeur. Une suite
croissante majorée par $10$ ne converge pas nécessairement vers $10$.
\end{avertissement}

\begin{theoreme}[limite d'une suite récurrente]
Si $u_{n+1} = g(u_{n})$, si $g$ est continue sur un intervalle contenant tous
les termes, et si $(u_{n})$ converge vers $\ell$, alors $g(\ell) = \ell$.
\end{theoreme}

\begin{demonstration}
La suite $(u_{n+1})$ converge vers $\ell$, comme $(u_{n})$. Par continuité de
$g$, $\left(g(u_{n})\right)$ converge vers $g(\ell)$. Or ces deux suites sont
égales terme à terme : leurs limites coïncident.
\end{demonstration}

\begin{methode}{encadrer l'écart à la limite}
Le schéma du problème d'analyse, toujours le même.
\begin{enumerate}[label=\textbf{\arabic*.}, leftmargin=*, itemsep=0.15em]
  \item Montrer que $g(x) = x$ admet une solution unique $\alpha$ sur $I$.
  \item Montrer par récurrence que tous les termes restent dans $I$.
  \item Majorer $\abs{g'}$ par $k<1$ sur $I$, puis appliquer l'inégalité des
        accroissements finis : $\abs{u_{n+1}-\alpha} \leq k\abs{u_{n}-\alpha}$.
  \item En déduire par récurrence $\abs{u_{n}-\alpha} \leq k^{n}\abs{u_{0}-\alpha}$,
        puis conclure à la convergence.
\end{enumerate}
\end{methode}

\begin{exemple}[la même suite, deux visages]
Soit $u_{0} = 5$ et $u_{n+1} = \dfrac{u_{n}+3}{4}$. Posons $v_{n} = u_{n}-1$ :
\[
  v_{n+1} = \frac{u_{n}+3}{4}-1 = \frac{u_{n}-1}{4} = \frac14 v_{n}.
\]
$(v_{n})$ est géométrique de raison $\frac14$ et de premier terme $v_{0} = 4$,
d'où $u_{n} = 1+4^{1-n}$, qui tend vers $1$.
\end{exemple}

\begin{entrainement}[suites et convergence]
\exo[\niveau{1}] Étudier le sens de variation de $u_{n} = \dfrac{n}{n+1}$ et
montrer qu'elle est bornée.

\exo[\niveau{2}] Soit $u_{0} = 2$ et $u_{n+1} = \frac12\left(u_{n}
+\frac{2}{u_{n}}\right)$. Montrer que $u_{n} \geq \sqrt2$ pour tout $n$, puis
que la suite décroît.
\end{entrainement}

\begin{corrige}
\textbf{1.} $u_{n+1}-u_{n} = \dfrac{1}{(n+1)(n+2)} > 0$ : croissante. Comme
$0 \leq u_{n} < 1$, elle est bornée.

\textbf{2.} $u_{n+1}-\sqrt2 = \dfrac{\left(u_{n}-\sqrt2\right)^{2}}{2u_{n}} \geq 0$
dès que $u_{n}>0$ : par récurrence, $u_{n} \geq \sqrt2$ pour tout $n$. Alors
$u_{n+1}-u_{n} = \dfrac{2-u_{n}^{2}}{2u_{n}} \leq 0$ puisque
$u_{n}^{2} \geq 2$ : la suite décroît.
\end{corrige}

% ===========================================================================
\section{Suites arithmétiques, géométriques et arithmético-géométriques}

\begin{definition}[suites arithmétiques et géométriques]
$(u_{n})$ est \textbf{arithmétique} de raison $r$ si $u_{n+1} = u_{n}+r$ pour
tout $n$ ; \textbf{géométrique} de raison $q$ si $u_{n+1} = q\,u_{n}$.
\end{definition}

\begin{propriete}[terme général et sommes]
Pour une suite arithmétique de premier terme $u_{0}$ :
\[
  u_{n} = u_{0}+nr, \qquad
  \sum_{k=0}^{n}u_{k} = (n+1)\,\frac{u_{0}+u_{n}}{2}.
\]
Pour une suite géométrique de raison $q \neq 1$ :
\[
  u_{n} = u_{0}\,q^{n}, \qquad
  \sum_{k=0}^{n}u_{k} = u_{0}\,\frac{1-q^{n+1}}{1-q}.
\]
\end{propriete}

\begin{propriete}[limite d'une suite géométrique]
$\limn q^{n} = 0$ si $\abs q<1$ ; $=+\infty$ si $q>1$ ; $=1$ si $q = 1$ ; et
la suite n'a pas de limite si $q \leq -1$.
\end{propriete}

\begin{definition}[suite arithmético-géométrique]
Une suite est \textbf{arithmético-géométrique} lorsqu'elle vérifie
$u_{n+1} = a\,u_{n}+b$, avec $a \neq 1$.
\end{definition}

\begin{propriete}[réduction]
Soit $\alpha$ le point fixe, solution de $\alpha = a\alpha+b$. Alors la suite
$v_{n} = u_{n}-\alpha$ est géométrique de raison $a$, et
$u_{n} = \alpha + \left(u_{0}-\alpha\right)a^{n}$.
\end{propriete}

\begin{remarque}
C'est exactement ce schéma qui décrira, au chapitre des mathématiques
financières, le capital restant dû sur un emprunt remboursé par mensualités
constantes.
\end{remarque}

\begin{entrainement}[arithmético-géométriques]
\exo[\niveau{1}] $u_{0} = 0$ et $u_{n+1} = 3u_{n}+4$ : déterminer $u_{n}$.

\exo[\niveau{2}] $u_{0} = 10$ et $u_{n+1} = 0{,}8u_{n}+30$ : déterminer la
limite et l'interpréter.

\exo[\niveau{3}] Une coopérative emprunte et rembourse chaque mois une somme
constante ; le capital restant dû vérifie $C_{n+1} = 1{,}02\,C_{n}-50\,000$,
avec $C_{0} = 2\,000\,000$. Déterminer $C_{n}$ en fonction de $n$.
\end{entrainement}

\begin{corrige}
\textbf{1.} Point fixe $\alpha = -2$ ; $u_{n} = -2+2\times3^{n}$.

\textbf{2.} Point fixe $\alpha = 150$ ; $u_{n}-150 = -140\times0{,}8^{n} \to 0$ :
la suite converge vers $150$.

\textbf{3.} Point fixe $\alpha = 2\,500\,000$ ;
$C_{n} = 2\,500\,000-500\,000\times1{,}02^{n}$.
\end{corrige}

% ===========================================================================
\section{Suites stationnaires et suites adjacentes}

\begin{definition}[suite stationnaire, suites adjacentes]
Une suite est \textbf{stationnaire} si elle est constante à partir d'un
certain rang. Deux suites $(u_{n})$ et $(v_{n})$ sont \textbf{adjacentes} si
l'une croît, l'autre décroît, et si $\limn\left(v_{n}-u_{n}\right) = 0$.
\end{definition}

\begin{theoreme}[théorème des suites adjacentes]
Deux suites adjacentes convergent vers la même limite $\ell$, et pour tout
$n$, $u_{n} \leq \ell \leq v_{n}$.
\end{theoreme}

\begin{entrainement}[suites adjacentes]
\exo[\niveau{2}] Montrer que $u_{n} = 1-\frac1n$ et $v_{n} = 1+\frac{1}{n^{2}}$
sont adjacentes.
\end{entrainement}

\begin{corrige}
$(u_{n})$ croît, $(v_{n})$ décroît, et $v_{n}-u_{n} = \frac1n+\frac{1}{n^{2}} \to 0$.
\end{corrige}

% ===========================================================================
\begin{annales}
La suite arrive en fin de problème d'analyse, une fois la fonction étudiée.
Les énoncés qui suivent sont repris de sujets réellement posés ; leur
provenance est indiquée à chaque fois.

\exoann{Baccalauréat, Congo-Brazzaville, série C, session 2022 — exercice 3}
\emph{On admet que la fonction $f$ du sujet envoie $\intervalle{2}{2{,}5}$
dans lui-même, que $-0{,}5 \leq f'(x) < 0$ sur cet intervalle, et que $\alpha$
en est l'unique point fixe.}

Soit $(u_{n})$ définie par $u_{0} = 2$ et $u_{n+1} = f(u_{n})$.
\begin{enumerate}[label=\textbf{\alph*)}, leftmargin=*, itemsep=0.15em]
  \item Montrer que si $u_{n} \in \intervalle{2}{2{,}5}$ alors
        $u_{n+1} \in \intervalle{2}{2{,}5}$.
  \item Montrer que $\abs{u_{n+1}-\alpha} \leq \frac12\abs{u_{n}-\alpha}$,
        puis que $\abs{u_{n}-\alpha} \leq
        \left(\frac12\right)^{n+1}$.
  \item En déduire la convergence, puis trouver $n_{0}$ tel que
        $\abs{u_{n}-\alpha} \leq 10^{-3}$.
\end{enumerate}

\exoann{Baccalauréat, Madagascar, série A, session 2026 — exercice 1}
Soit $(U_{n})$ définie par $U_{0} = 5$ et
$U_{n+1} = \dfrac{U_{n}+3}{4}$.
\begin{enumerate}[label=\textbf{\alph*)}, leftmargin=*, itemsep=0.15em]
  \item On pose $V_{n} = U_{n}-1$. Montrer que $(V_{n})$ est géométrique et
        préciser sa raison.
  \item Exprimer $V_{n}$ puis $U_{n}$ en fonction de $n$, et donner la limite
        de $(U_{n})$.
  \item Calculer $S_{n} = U_{0}+U_{1}+\cdots+U_{n}$.
\end{enumerate}

\exoann{Baccalauréat, Burkina Faso, série D, session 2023 — exercice 2}
Soit $(U_{n})$ définie par $U_{0} = 2$ et
$\ln\!\left(U_{n+1}\right) = 1+\ln\!\left(U_{n}\right)$.
\begin{enumerate}[label=\textbf{\alph*)}, leftmargin=*, itemsep=0.15em]
  \item Montrer que $U_{n+1} = eU_{n}$, puis que $(U_{n})$ est géométrique.
  \item Exprimer $U_{n}$ en fonction de $n$, étudier sa monotonie et sa
        limite.
  \item Calculer $S_{n} = U_{0}+\cdots+U_{n}$.
\end{enumerate}

\exoann{Baccalauréat, Côte d'Ivoire, série C, session 2024 — exercice 2}
Déterminer, en justifiant, la limite de la suite définie par
$u_{n} = \left(-\dfrac23\right)^{n}$.

\exoann[suite géométrique et seuil]{Baccalauréat, France, série ES, Antilles-Guyane, 22 juin 2016 — exercice 4}
Un département a contraint dès 2013 certaines entreprises à diminuer chaque
année la quantité de polluants qu'elles rejettent dans l'air. Ces
entreprises ont rejeté $410$ tonnes de polluants en 2013 et $332$ tonnes en
2015 ; on considère que le taux annuel de diminution est constant.
\begin{enumerate}[label=\textbf{\arabic*.}, leftmargin=*, itemsep=0.2em]
  \item Justifier que ce taux de diminution annuel peut être considéré comme
        égal à $10\,\%$.
  \item En admettant que ce taux de $10\,\%$ reste constant, déterminer à
        partir de quelle année la quantité rejetée ne dépassera plus le seuil
        de $180$ tonnes fixé par le conseil départemental.
\end{enumerate}

\exoann[suite arithmético-géométrique]{Baccalauréat, France, série ES/L, Liban, 27 mai 2014 — exercice 3, parties 1 et 2 (extrait)}
Une médiathèque a ouvert ses portes en 2013 et a enregistré $2\,500$
inscriptions cette année-là. Elle estime que, chaque année, $80\,\%$ des
anciens inscrits renouvelleront leur inscription l'année suivante et qu'il y
aura $400$ nouveaux adhérents. On note $a_{n}$ le nombre d'inscrits pendant
l'année $2013+n$, avec $a_{0} = 2\,500$.
\begin{enumerate}[label=\textbf{\arabic*.}, leftmargin=*, itemsep=0.2em]
  \item Calculer $a_{1}$ et $a_{2}$, puis justifier que, pour tout entier
        naturel $n$, $a_{n+1} = 0{,}8\,a_{n}+400$.
  \item On pose $v_{n} = a_{n}-2\,000$. Démontrer que $(v_{n})$ est
        géométrique de raison $0{,}8$, en déduire $a_{n}$ en fonction de $n$,
        puis la limite de $(a_{n})$.
\end{enumerate}
\end{annales}

\begin{corrige}[annales]
\textbf{1.} La stabilité de l'intervalle est donnée par l'énoncé. Comme
$\abs{f'} \leq \frac12$, l'inégalité des accroissements finis donne le
facteur $\frac12$, puis la récurrence
$\abs{u_{n}-\alpha} \leq \left(\frac12\right)^{n}\abs{u_{0}-\alpha}$ ; enfin
$\abs{u_{0}-\alpha} \leq \frac12$ puisque $u_{0}$ et $\alpha$ sont dans un
intervalle de longueur $\frac12$. La convergence suit, et
$\left(\frac12\right)^{n+1} \leq 10^{-3}$ donne $n \geq 9$.

\textbf{2.} $V_{n+1} = \frac{U_{n}-1}{4} = \frac14V_{n}$ : suite géométrique
de raison $\frac14$, de premier terme $V_{0} = 4$. Donc
$V_{n} = 4\left(\frac14\right)^{n}$ et $U_{n} = 1+4^{1-n}$, qui tend vers
$1$. Enfin
$S_{n} = (n+1)+\frac{16}{3}\left(1-4^{-(n+1)}\right)$.

\textbf{3.} $\ln U_{n+1} = \ln(eU_{n})$ donne $U_{n+1} = eU_{n}$ : suite
géométrique de raison $e>1$, donc $U_{n} = 2e^{n}$, croissante, de limite
$+\infty$. La somme vaut $S_{n} = 2\dfrac{e^{n+1}-1}{e-1}$.

\textbf{4.} $\abs{-\frac23} = \frac23<1$, donc la suite converge vers $0$ ;
elle est alternée, ses termes changeant de signe à chaque rang.

\textbf{5.} Sur deux ans, $\dfrac{332}{410} \approx 0{,}8098$. Or
$0{,}9^{2} = 0{,}81$, valeur très proche : on retient un taux de diminution
annuel de $10\,\%$, si bien que le rejet de l'année $2013+n$ vaut
$u_{n} = 410 \times 0{,}9^{n}$. L'inéquation $u_{n} \leq 180$ équivaut à
$0{,}9^{n} \leq \frac{180}{410} \approx 0{,}4390$, soit, en passant au
logarithme népérien (décroissant sur $\intervalleo0{+\infty}$ car
$0{,}9<1$),
\[
  n \geq \frac{\ln(0{,}4390)}{\ln(0{,}9)} \approx 7{,}82 .
\]
Le plus petit entier convenable est $n = 8$ : c'est en $2013+8 = 2021$ que le
rejet passe pour la première fois sous le seuil de $180$ tonnes
($u_{7}\approx196$ tonnes, encore au-dessus ; $u_{8}\approx176{,}5$ tonnes).

\textbf{6. 1.} $a_{1} = 0{,}8\times2\,500+400 = 2\,400$ et
$a_{2} = 0{,}8\times2\,400+400 = 2\,320$. La relation
$a_{n+1} = 0{,}8\,a_{n}+400$ traduit exactement l'énoncé : $80\,\%$ des
$a_{n}$ anciens inscrits renouvellent, et $400$ nouveaux s'ajoutent.
\textbf{2.} $v_{n+1} = a_{n+1}-2\,000 = 0{,}8\,a_{n}-1\,600
= 0{,}8\,(a_{n}-2\,000) = 0{,}8\,v_{n}$ : $(v_{n})$ est géométrique de raison
$0{,}8$ et de premier terme $v_{0} = 500$. Donc
$v_{n} = 500\times0{,}8^{n}$ et $a_{n} = 500\times0{,}8^{n}+2\,000$. Comme
$\abs{0{,}8}<1$, $0{,}8^{n}\to0$, donc $a_{n}\to2\,000$ : le nombre
d'inscrits se stabiliserait autour de $2\,000$ si le schéma se maintenait.
\end{corrige}

\begin{jecherche}[le remboursement d'un prêt agricole]
Une coopérative emprunte $12\,000\,000$ d'ariary à un taux mensuel de
$1\,\%$. Elle rembourse chaque mois une mensualité constante de $M$ ariary,
versée juste après que les intérêts du mois ont été ajoutés. On note $C_{n}$
le capital restant dû après le $n$-ième versement, et $C_{0} = 12\,000\,000$.

\begin{enumerate}[label=\textbf{\arabic*.}, leftmargin=*, itemsep=0.3em]
  \item Justifier que $C_{n+1} = 1{,}01\,C_{n}-M$.
  \item Déterminer le point fixe de cette relation et l'interpréter : que se
        passe-t-il si $M = 120\,000$ ?
  \item On prend $M = 300\,000$. Exprimer $C_{n}$ en fonction de $n$.
  \item Au bout de combien de mois le prêt est-il soldé ?
\end{enumerate}
\end{jecherche}

\begin{corrige}[le remboursement d'un prêt agricole]
\textbf{1.} Les intérêts du mois valent $0{,}01\,C_{n}$ ; on ajoute puis on
retire la mensualité : $C_{n+1} = C_{n}+0{,}01C_{n}-M = 1{,}01C_{n}-M$.

\textbf{2.} Le point fixe vérifie $\alpha = 1{,}01\alpha-M$, soit
$\alpha = 100M$. Pour $M = 120\,000$, $\alpha = 12\,000\,000 = C_{0}$ : le
capital reste constant, la mensualité ne paie que les intérêts et le prêt
n'est jamais remboursé.

\textbf{3.} Pour $M = 300\,000$, $\alpha = 30\,000\,000$ et
$C_{n} = 30\,000\,000-18\,000\,000\times1{,}01^{n}$.

\textbf{4.} $C_{n} \leq 0$ équivaut à $1{,}01^{n} \geq \frac53$, soit
$n \geq \frac{\ln(5/3)}{\ln1{,}01} \approx 51{,}3$ : le prêt est soldé au
cours du $52^{\text{e}}$ mois.
\end{corrige}

% ===========================================================================
\begin{jemeteste}
Répondre par vrai ou faux, en justifiant en une ligne.

\begin{enumerate}[label=\textbf{\arabic*.}, leftmargin=*, itemsep=0.28em]
  \item Toute suite croissante tend vers $+\infty$.
  \item Une suite croissante et majorée par $10$ converge vers $10$.
  \item Si $u_{n+1} = g(u_{n})$ et si $(u_{n})$ converge vers $\ell$ avec $g$
        continue, alors $g(\ell) = \ell$.
  \item Une suite bornée converge.
  \item Deux suites adjacentes ont la même limite.
\end{enumerate}
\end{jemeteste}

\begin{corrige}[je me teste]
\textbf{1.} Faux — elle peut être majorée, et converger.
\textbf{2.} Faux — elle converge vers un réel inférieur ou égal à $10$.
\textbf{3.} Vrai — c'est le théorème du point fixe.
\textbf{4.} Faux — $\left((-1)^{n}\right)$ est bornée sans converger.
\textbf{5.} Vrai — c'est la conclusion du théorème des suites adjacentes.
\end{corrige}

% ===========================================================================
\begin{commeaubac}
\bmtsource{Exercice au format de l'épreuve — série OSE}
Soit $f$ la fonction du problème, continue sur $\intervalle{1}{2}$, à valeurs
dans $\intervalle{1}{2}$, et telle que $\abs{f'(x)} \leq \frac23$ sur cet
intervalle. On pose $h(x) = f(x)-x$.

\begin{enumerate}[label=\textbf{\arabic*.}, leftmargin=*, itemsep=0.25em]
  \item À l'aide de l'étude de $h$ sur $\intervalle{1}{2}$, montrer que
        l'équation $h(x) = 0$ admet une solution unique
        $\alpha \in \intervalle{1}{2}$. \hfill\textup{(1 pt)}
  \item Soit $(U_{n})$ définie par $U_{0} = 1$ et $U_{n+1} = f(U_{n})$.
        \begin{enumerate}[label=\textbf{\alph*)}, leftmargin=1.4em]
          \item Démontrer par récurrence que $U_{n} \in \intervalle{1}{2}$.
                \hfill\textup{(0,75 pt)}
          \item Démontrer que
                $\abs{U_{n+1}-\alpha} \leq \frac23\abs{U_{n}-\alpha}$, puis
                que $\abs{U_{n}-\alpha} \leq \left(\frac23\right)^{n}$.
                \hfill\textup{(1 pt)}
          \item En déduire l'étude de la convergence de $(U_{n})$.
                \hfill\textup{(0,5 pt)}
          \item Déterminer le plus petit entier $n_{0}$ tel que $U_{n_{0}}$
                soit une valeur approchée de $\alpha$ à $10^{-3}$ près.
                \hfill\textup{(0,75 pt)}
        \end{enumerate}
\end{enumerate}
\end{commeaubac}

\begin{corrige}[exercice au format de l'épreuve]
\textbf{1.} $h$ est continue sur $\intervalle{1}{2}$ ; comme $f$ y prend ses
valeurs, $h(1) = f(1)-1 \geq 0$ et $h(2) = f(2)-2 \leq 0$. De plus
$h'(x) = f'(x)-1 \leq \frac23-1<0$ : $h$ est strictement décroissante. Le
théorème des valeurs intermédiaires donne l'existence, la stricte monotonie
l'unicité.

\textbf{2. a)} $U_{0} = 1$ appartient à l'intervalle ; si $U_{n}$ y
appartient, $U_{n+1} = f(U_{n})$ aussi puisque $f$ y prend ses valeurs.
\textbf{b)} L'inégalité des accroissements finis appliquée entre $U_{n}$ et
$\alpha$ donne le facteur $\frac23$ ; la récurrence donne ensuite
$\abs{U_{n}-\alpha} \leq \left(\frac23\right)^{n}\abs{U_{0}-\alpha}$, et
$\abs{U_{0}-\alpha} \leq 1$.
\textbf{c)} Le majorant tend vers $0$ : la suite converge vers $\alpha$.
\textbf{d)} $\left(\frac23\right)^{n} \leq 10^{-3}$ donne
$n \geq \frac{3\ln10}{\ln1{,}5} \approx 17{,}03$, soit $n_{0} = 18$.
\end{corrige}
